\section{Evaluation}
\label{sec:evaluation}

Our approach is implemented as a language called \OurLang, implemented in PureScript. All benchmarks were run
on a Dell Latitude 5411, with an Intel Core i7-10850H at 2.70GHz and 16Gb of RAM. We ran our benchmarks on
Google Chome 118.0.5993.117, running Javascript V8 11.8.172.16.

We assess the benefits of our solution by answering the following three research questions:
\begin{itemize}
    \item RQ1: How much does slicing performance improve in the graphical approach?
    \item RQ2: How much overhead does the construction of the graph introduce?
    \item RQ3: How do the equivalent slicing implementations compare to each other?
\end{itemize}

\subsection{Benchmarks} \label{sec:benchmarks}
We have a selection of 67 benchmarks, ranging from toy programs to more complex
algorithms from computer science literature, \todo{which we manually translated to Fluid?}.
%Since we have written our own language, we needed to construct our own benchmarks.
All the times presented in this
section are measured in milliseconds. We present the full suite of benchmarks in
Table \ref{table:all-benches}.
\todo{Need to further describe the range of benchmarks, and what they represent}

\subsection{RQ1: How much does slicing performance improve in the graphical approach?}
%\subsection{Slicing Performance}

Compared to the Galois slicing approach of \citet{perera22} based on running adjoint analyses over an
execution trace, our approach factors the analyis into two components: an operational semantics that evaluates
a program to a DDG, and a pair of slicing operators for computing reachability and its adjoint over the graph.
We show here that although the overhead of producing a DDG is comparable to that of producing a trace, the
slicing algorithms run on average much faster than their trace-based counterparts, making them suitable for
interactive applcations like the related inputs and related outputs queries discussed in \secref{example}. We
have already established in \secref{graph} that slicing in the graphical setting is much simplified
algorithmically. Since we are targeting interactive systems, we therefore feel the gain in slicing speed is
worth the loss of performance in evaluation.

\todo{performance in evaluation?}

Our metric for the improvement to slicing speed is based on a well-known set of
criteria proposed in \cite{nielsen93}, which we paraphrase here:

\begin{itemize}
    \item \textbf{0.1 seconds} or less feels instantaneous to the user
    \item \textbf{1.0 seconds} is the limit before the users flow of thought feels
    interrupted
    \item \textbf{10.0 seconds} is the limit for keeping the users attention. Anything longer
    than this and the user will shift their focus to other tasks
\end{itemize}

As we report results in measurements of milliseconds, then slicing benchmark results
of 100 or less constitute near-instantaneous interaction with the system.

In Table \ref{table:slicing-speedups}, we see that in general we do reach
the performance guideline specified by Nielsen. Even in our worst cases - the
benchmarks prefixed with \textit{graphics/} - we do not ever tip over the one
second mark. The performance on those cases is still worse than other benchmarks
but still within an acceptable range.

One important point of discussion is the range in slicing speedups we observe.
\todo{This description of some of the benchmarks needs to go to sec~\ref{sec:benchmarks}.}
For example, in the cases prefixed \textit{slicing/convolution} we implement
some classic convolutional filters from computer vision literature, an edge
detection filter, and an embossing filter. We see that these cases produce the
most impressive speedup numbers, with forward slicing speeding up by between
160 and 380 times, and backward slicing speeding up between 100 and 120 times.

Contrast this with the benchmark case \textit{slicing/dtw/compute-dtw}, the
example we presented earlier in the paper. This is one of the largest programs
we have tested, and we see that the speedup in both forward and backward slicing
is approximately a factor of eight.

We believe that the disparity in performance improvements is due to an inherent
data-parallelism of certain examples. In an image convolution, there are a large
number of cells in an output image which do not share any cells of the input.

Contrast this with the Dynamic Time Warp algorithm, which is inherently more
sequential. We therefore propose that our method is more suited to the efficient
program slicing of algorithms - or steps of algorithms - which exhibit a high
degree of data-independence.
\todo{sure if this should come before or
after the bit where we discuss the competing implementations we need.}

{\tiny\begin{longtable}{lrrrrrr}
\caption{Slicing times and average speedups} \label{table:slicing-speedups} \\
\toprule
 & Trace-Fwd & Graph-Fwd & Fwd-Speedup & Trace-Bwd & Graph-Bwd & Bwd-Speedup \\
Test-Name &  &  &  &  &  &  \\
\midrule
\endfirsthead
\caption[]{Slicing times and average speedups} \\
\toprule
 & Trace-Fwd & Graph-Fwd & Fwd-Speedup & Trace-Bwd & Graph-Bwd & Bwd-Speedup \\
Test-Name &  &  &  &  &  &  \\
\midrule
\endhead
\midrule
\multicolumn{7}{r}{Continued on next page} \\
\midrule
\endfoot
\bottomrule
\endlastfoot
arithmetic & 0.57 & 0.17 & 3.40 & 4.83 & 0.77 & 6.30 \\
array & 1.10 & 0.03 & 33.00 & 3.87 & 0.20 & 19.33 \\
compose & 0.20 & 0.00 & inf & 3.43 & 0.30 & 11.44 \\
dicts & 0.50 & 0.60 & 0.83 & 3.97 & 0.30 & 13.22 \\
div-mod-quot-rem & 1.53 & 0.67 & 2.30 & 3.03 & 0.30 & 10.11 \\
factorial & 2.20 & 0.00 & inf & 3.70 & 0.23 & 15.86 \\
filter & 1.47 & 0.07 & 22.00 & 3.93 & 0.17 & 23.60 \\
first-class-constr & 1.63 & 0.10 & 16.33 & 3.23 & 0.63 & 5.11 \\
flatten & 2.83 & 0.20 & 14.17 & 3.83 & 0.40 & 9.58 \\
Foldr-sumSquares & 2.20 & 0.00 & inf & 4.30 & 0.23 & 18.43 \\
lexicalScoping & 0.30 & 0.00 & inf & 2.50 & 0.57 & 4.41 \\
length & 0.53 & 0.03 & 16.00 & 3.03 & 0.27 & 11.37 \\
lookup & 1.90 & 0.03 & 57.00 & 2.77 & 0.20 & 13.83 \\
map & 1.93 & 0.07 & 29.00 & 4.07 & 0.33 & 12.20 \\
mergeSort & 9.90 & 0.07 & 148.50 & 5.07 & 0.23 & 21.71 \\
normalise & 0.43 & 0.10 & 4.33 & 3.87 & 0.20 & 19.33 \\
pattern-match & 2.23 & 0.07 & 33.50 & 3.43 & 0.23 & 14.71 \\
range & 7.60 & 0.13 & 57.00 & 8.17 & 0.23 & 35.00 \\
records & 2.03 & 0.07 & 30.50 & 3.60 & 0.20 & 18.00 \\
reverse & 0.53 & 0.00 & inf & 2.30 & 0.17 & 13.80 \\
desugar/list-comp-1 & 11.07 & 0.27 & 41.50 & 12.13 & 0.37 & 33.09 \\
desugar/list-comp-2 & 27.87 & 0.80 & 34.83 & 26.60 & 0.37 & 72.55 \\
desugar/list-comp-3 & 6.07 & 0.00 & inf & 7.10 & 0.23 & 30.43 \\
desugar/list-comp-4 & 3.17 & 0.03 & 95.00 & 4.50 & 0.43 & 10.38 \\
desugar/list-comp-5 & 1.60 & 0.07 & 24.00 & 3.27 & 0.13 & 24.50 \\
desugar/list-comp-6 & 0.17 & 0.03 & 5.00 & 3.27 & 0.27 & 12.25 \\
desugar/list-comp-7 & 2.70 & 0.00 & inf & 3.50 & 0.23 & 15.00 \\
desugar/list-enum & 1.33 & 0.07 & 20.00 & 3.00 & 0.27 & 11.25 \\
slicing/add & 0.07 & 0.07 & 1.00 & 2.67 & 0.33 & 8.00 \\
slicing/array/lookup & 10.77 & 0.03 & 323.00 & 9.30 & 0.43 & 21.46 \\
slicing/array/dims & 0.90 & 0.17 & 5.40 & 2.60 & 0.47 & 5.57 \\
slicing/convolution/edgeDetect & 685.13 & 4.03 & 169.87 & 546.50 & 5.33 & 102.47 \\
slicing/convolution/emboss & 628.10 & 1.90 & 330.58 & 439.73 & 3.70 & 118.85 \\
slicing/convolution/gaussian & 602.27 & 1.60 & 376.42 & 432.80 & 4.00 & 108.20 \\
slicing/dict/create & 0.00 & 0.07 & 0.00 & 2.80 & 0.27 & 10.50 \\
slicing/dict/difference & 0.23 & 0.10 & 2.33 & 2.67 & 0.40 & 6.67 \\
slicing/dict/disjointUnion & 0.07 & 0.07 & 1.00 & 1.77 & 0.23 & 7.57 \\
slicing/dict/foldl & 2.97 & 0.03 & 89.00 & 4.90 & 0.27 & 18.37 \\
slicing/dict/intersectionWith & 0.33 & 0.20 & 1.67 & 3.10 & 0.30 & 10.33 \\
slicing/dict/fromRecord & 0.03 & 0.10 & 0.33 & 1.77 & 0.27 & 6.63 \\
slicing/dict/get & 0.20 & 0.03 & 6.00 & 2.50 & 0.23 & 10.71 \\
slicing/dict/map & 0.67 & 0.07 & 10.00 & 2.60 & 0.43 & 6.00 \\
slicing/divide & 0.03 & 0.10 & 0.33 & 2.70 & 0.20 & 13.50 \\
slicing/filter & 1.77 & 0.30 & 5.89 & 2.63 & 0.27 & 9.88 \\
slicing/intersperse & 1.07 & 0.20 & 5.33 & 1.97 & 0.27 & 7.38 \\
slicing/intersperse & 1.20 & 0.27 & 4.50 & 2.53 & 0.43 & 5.85 \\
slicing/length & 0.57 & 0.50 & 1.13 & 2.90 & 0.43 & 6.69 \\
slicing/list-comp & 11.87 & 1.10 & 10.79 & 11.17 & 0.47 & 23.93 \\
slicing/list-comp & 10.73 & 0.97 & 11.10 & 9.83 & 0.43 & 22.69 \\
slicing/lookup & 1.83 & 0.33 & 5.50 & 2.23 & 0.47 & 4.79 \\
slicing/map & 0.37 & 0.13 & 2.75 & 2.37 & 0.37 & 6.45 \\
slicing/multiply & 0.03 & 0.00 & inf & 2.10 & 0.23 & 9.00 \\
slicing/nth & 0.50 & 0.07 & 7.50 & 2.50 & 0.17 & 15.00 \\
slicing/section-5-example & 12.10 & 0.93 & 12.96 & 9.97 & 0.57 & 17.59 \\
slicing/section-5-example & 12.50 & 0.47 & 26.79 & 11.87 & 0.47 & 25.43 \\
slicing/section-5-example & 12.10 & 0.60 & 20.17 & 10.10 & 0.33 & 30.30 \\
slicing/zeros & 0.47 & 0.43 & 1.08 & 2.93 & 0.40 & 7.33 \\
slicing/zeros & 0.50 & 0.17 & 3.00 & 2.33 & 0.33 & 7.00 \\
slicing/zipWith & 1.87 & 0.40 & 4.67 & 3.00 & 0.37 & 8.18 \\
slicing/matrix-update & 36.23 & 0.33 & 108.70 & 30.20 & 0.47 & 64.71 \\
slicing/dtw/compute-dtw & 154.23 & 17.47 & 8.83 & 145.00 & 18.53 & 7.82 \\
graphics/background & 1.30 & 0.63 & 2.05 & 2.93 & 1.17 & 2.51 \\
graphics/grouped-bar-chart & 81.97 & 142.10 & 0.58 & 64.63 & 1.73 & 37.29 \\
graphics/line-chart & 112.37 & 614.03 & 0.18 & 84.80 & 1.33 & 63.60 \\
graphics/stacked-bar-chart & 48.43 & 103.87 & 0.47 & 36.73 & 1.13 & 32.41 \\
\end{longtable}
}

\subsection{RQ2: How much overhead does the construction of the graph introduce?}
%\subsection{Overhead}

\todo{This is confusing because the RQ talks about graph construction, whereas in the actual section, you talk about evaluation (this also relates to one of my comments above).}

In both the trace-based and graphical approaches, we construct an auxiliary data
structure during evaluation. In the trace-based approach, we construct what is
in essence a simplified version of the AST of the program, only representing
things such as the branch of a pattern match or conditional expression which
was taken.

In contrast, the graphical evaluation approach works by annotating values with
vertices, and building sets of neighbors. We collect the sets of vertices and
edges produced at the end of evaluation, and only afterward construct the actual
representation of the graph.

We can see in Table \ref{table:eval-speedups} that across our test-cases, trace
based evaluation takes about $54\%$ of the time that graph evaluation takes. We
believe that this is an acceptable level of performance loss, given the gains we
make in slicing performance, which we discuss next.
\todo{I need to move/reformat the table because it currently takes up a whole page}
{\tiny\begin{longtable}{lrrr}
\caption{Evaluation Time of Traces Versus Graphs (ms)} \label{table:eval-speedups} \\
\toprule
 & T-Eval & G-Eval & Eval-Speedup \\
\midrule
\endfirsthead
\caption[]{Evaluation Time of Traces Versus Graphs (ms)} \\
\toprule
 & T-Eval & G-Eval & Eval-Speedup \\
\midrule
\endhead
\midrule
\multicolumn{4}{r}{Continued on next page} \\
\midrule
\endfoot
\bottomrule
\endlastfoot
slicing/dtw/compute-dtw & 65.03 & 219.40 & 0.30 \\
slicing/convolution/edgeDetect & 330.10 & 1039.63 & 0.32 \\
slicing/convolution/emboss & 290.47 & 849.30 & 0.34 \\
slicing/convolution/gaussian & 293.93 & 866.70 & 0.34 \\
Matrix-Avg & 244.88 & 743.76 & 0.33 \\
graphics/background & 0.27 & 2.30 & 0.12 \\
graphics/grouped-bar-chart & 39.97 & 912.83 & 0.04 \\
graphics/line-chart & 59.57 & 554.17 & 0.11 \\
graphics/stacked-bar-chart & 23.93 & 378.23 & 0.06 \\
Graphics-Avg & 30.93 & 461.88 & 0.07 \\
\end{longtable}
}

\subsection{RQ3: How do the equivalent slicing implementations compare to each other?}
%\subsection{Equivalent Implementations}
As we have discussed earlier, if one is careful about the definition of inputs
and outputs on a given Dynamic Dependence Graph, then we have multiple ways to
implement equivalent operators. As part of our testing, we compared equivalent
implementations.

First, we considered implementing forward slicing as the De Morgan dual of backward
slicing on the opposite graph. This is identified in our benchmarks as \NaiveFwd.
We also implemented the more direct algorithm described in \secref{graph}, which
is what we marked as \GraphFwd. We also implement twin implementations of the
duals of slicing. One (pair of) implementations, that of complementing, running the slicing
operation, then complementing again, is denoted as \GraphBwdCmp, and \GraphFwdCmp,
for the duals ofbackward and forward slicing respectively. The implementations
using the opposite slicing procedure on the opposite graph are denoted as
\GraphBwdOp, and \GraphFwdOp.

In our implementation, the operation for taking a graph's opposite is much
cheaper than that of complementation. This is reflected in our performance
results where we see that the implementation using the opposite graph is much
quicker generally than the implementation using complements. Results for these
comparisons are contained in \tableref{equivalent-impls}.

{\tiny\begin{longtable}{lrrrr}
\caption{Equivalent implementations of Demanded By} \label{table:equivalent-impls} \\
\toprule
 & G-DemandedBy & G-Suff-Dual & T-DemandedBy & Naive-Fwd \\
\midrule
\endfirsthead
\caption[]{Equivalent implementations of Demanded By} \\
\toprule
 & G-DemandedBy & G-Suff-Dual & T-DemandedBy & Naive-Fwd \\
\midrule
\endhead
\midrule
\multicolumn{5}{r}{Continued on next page} \\
\midrule
\endfoot
\bottomrule
\endlastfoot
slicing/dtw/compute-dtw & 19203.07 & 108.97 & 179.67 & 156.00 \\
slicing/convolution/edgeDetect & 485.03 & 1156.10 & 1091.67 & 2870.77 \\
slicing/convolution/emboss & 440.37 & 853.77 & 719.90 & 2504.03 \\
slicing/convolution/gaussian & 440.17 & 854.07 & 905.40 & 2462.97 \\
Matrix-Avg & 5142.16 & 743.23 & 724.16 & 1998.44 \\
graphics/background & 0.07 & 1.40 & 0.23 & 1.30 \\
graphics/grouped-bar-chart & 2.87 & 1059.80 & 39.47 & 844.27 \\
graphics/line-chart & 1.57 & 450.30 & 52.30 & 702.07 \\
graphics/stacked-bar-chart & 1.67 & 600.27 & 22.40 & 331.17 \\
Graphics-Avg & 1.54 & 527.94 & 28.60 & 469.70 \\
\end{longtable}
}

\subsection{Discussion}
\todo{I think we still want to include some finishing discussion, but unsure if
this is appropriate. Could be a conclusion?}
We can see that whilst the graphical slicing paradigm can improve the performance
of program slicing dramatically, the increase in performance is varied. It
initially appeared that speedup in slicing time was inherently tied to the
size of the dynamic dependence graph and equivalent trace, but this is not true.

Instead the computational speedup is more closely tied to the structural
properties of the program of interest, and it's resultant dependence graph
We have tested on several non-toy programs, and we note that there is a
high variance in speedup of both forward and backward slicing time. We see that
some examples - such as Dynamic Time Warping - experience approximately an
eighttimes performance increase in both forward and backward slicing, whereas
our convolution examples demonstrate speed increases on the order of 100s.

The reason for this disparity is the inherent data-parallelism of certain
computations. A single iteration of a convolution creates very "local" dependency
structure. A cell in the top left of the output matrix only depends on a small
section of the input image. We therefore expect that the graph produced
in evaluation has lots of inputs which do not cooperate.

In contrast, the Dynamic Time Warping algorithms output matching
requires traversing the entire cost matrix computed in evaluation, and a cell in
the cost matrix relies on the computation of it's left, upward, and diagonal
neighbors. The implication of this is that the data depdendency through the
the computation is in a sense more sequential than that of, for example
convolution.
{\tiny\todo{Need to move this to appendix}
\begin{table}
\caption{All test-cases and all benchmarks, TODO: needs formatting somehow}
\label{table:all-benches}
\begin{tabular}{lrrrrrrrrrr}
\toprule
 & G-DemandedBy-Dir & G-DemandedBy-Suff & G-Demands & G-Eval & G-Suffices & Graph-Nodes & T-DemandedBy & T-Demands & T-Eval & T-Suffices \\
Test-Name &  &  &  &  &  &  &  &  &  &  \\
\midrule
desugar/list-comp-1 & 0.50 & 5.10 & 0.30 & 16.80 & 0.57 & 268.00 & 5.67 & 6.50 & 7.30 & 5.60 \\
desugar/list-comp-2 & 0.27 & 11.70 & 0.20 & 38.07 & 0.57 & 835.00 & 11.50 & 12.20 & 12.57 & 11.03 \\
desugar/list-comp-3 & 0.17 & 3.33 & 0.07 & 7.90 & 0.37 & 260.00 & 2.47 & 2.43 & 2.50 & 2.53 \\
desugar/list-comp-4 & 0.10 & 1.37 & 0.10 & 3.27 & 0.30 & 111.00 & 1.03 & 0.90 & 1.03 & 0.97 \\
desugar/list-comp-5 & 0.13 & 1.33 & 0.03 & 5.50 & 0.40 & 80.00 & 2.03 & 1.30 & 1.50 & 1.87 \\
desugar/list-comp-6 & 0.00 & 0.10 & 0.00 & 0.17 & 0.03 & 8.00 & 0.00 & 0.10 & 0.03 & 0.03 \\
desugar/list-comp-7 & 0.17 & 1.23 & 0.07 & 2.53 & 0.37 & 93.00 & 0.63 & 0.70 & 0.73 & 0.83 \\
desugar/list-enum & 0.07 & 0.67 & 0.07 & 2.13 & 0.10 & 63.00 & 0.53 & 0.47 & 0.53 & 0.83 \\
arithmetic & 0.03 & 0.17 & 0.07 & 0.53 & 0.07 & 23.00 & 0.10 & 0.17 & 0.10 & 0.07 \\
array & 0.10 & 0.33 & 0.00 & 0.97 & 0.13 & 41.00 & 0.13 & 0.40 & 0.30 & 0.17 \\
compose & 0.03 & 0.13 & 0.03 & 0.27 & 0.10 & 13.00 & 0.03 & 0.07 & 0.03 & 0.17 \\
dicts & 0.03 & 0.60 & 0.10 & 1.13 & 0.13 & 59.00 & 0.20 & 0.43 & 0.27 & 0.10 \\
div-mod-quot-rem & 0.10 & 2.27 & 0.03 & 2.80 & 0.73 & 150.00 & 0.23 & 0.20 & 0.33 & 0.33 \\
factorial & 0.07 & 0.87 & 0.07 & 2.07 & 0.03 & 97.00 & 0.73 & 0.57 & 0.97 & 0.60 \\
filter & 0.00 & 0.53 & 0.07 & 1.43 & 0.13 & 48.00 & 0.30 & 0.33 & 0.40 & 0.33 \\
first-class-constr & 0.10 & 0.70 & 0.03 & 1.73 & 0.13 & 66.00 & 0.27 & 0.27 & 0.43 & 0.50 \\
flatten & 0.07 & 1.13 & 0.00 & 2.30 & 0.13 & 87.00 & 0.67 & 0.57 & 0.70 & 0.73 \\
foldr-sumSquares & 0.07 & 0.53 & 0.03 & 1.30 & 0.10 & 57.00 & 0.47 & 0.43 & 0.33 & 0.40 \\
include-input-into-output & 0.07 & 0.03 & 0.00 & 0.10 & 0.03 & 3.00 & 0.00 & 0.00 & 0.00 & 0.03 \\
lexicalScoping & 0.03 & 0.10 & 0.00 & 0.33 & 0.10 & 14.00 & 0.00 & 0.10 & 0.17 & 0.07 \\
length & 0.00 & 0.27 & 0.00 & 0.40 & 0.03 & 24.00 & 0.07 & 0.10 & 0.17 & 0.30 \\
lookup & 0.17 & 1.40 & 0.07 & 2.13 & 0.27 & 96.00 & 0.50 & 0.53 & 0.60 & 0.70 \\
map & 0.07 & 0.60 & 0.10 & 1.43 & 0.10 & 69.00 & 0.33 & 0.43 & 0.40 & 0.53 \\
mergeSort & 0.10 & 1.80 & 0.07 & 6.07 & 0.40 & 160.00 & 2.47 & 1.33 & 2.67 & 2.80 \\
normalise & 0.03 & 0.27 & 0.07 & 0.57 & 0.03 & 28.00 & 0.10 & 0.17 & 0.13 & 0.17 \\
nub & 0.07 & 1.50 & 0.07 & 4.37 & 0.20 & 128.00 & 1.50 & 1.17 & 1.53 & 1.53 \\
pattern-match & 0.07 & 2.17 & 0.03 & 2.63 & 0.63 & 105.00 & 0.70 & 0.47 & 0.67 & 0.67 \\
range & 0.17 & 2.63 & 0.07 & 7.20 & 0.40 & 210.00 & 2.23 & 2.13 & 2.40 & 2.40 \\
records & 0.20 & 1.80 & 0.10 & 2.90 & 0.53 & 115.00 & 0.67 & 0.80 & 0.73 & 0.60 \\
record-lookup & 0.07 & 0.33 & 0.03 & 0.50 & 0.10 & 29.00 & 0.13 & 0.13 & 0.20 & 0.23 \\
reverse & 0.03 & 0.40 & 0.03 & 0.90 & 0.07 & 33.00 & 0.20 & 0.23 & 0.30 & 0.20 \\
slicing/add & 0.03 & 0.03 & 0.10 & 0.10 & 0.10 & 11.00 & 0.00 & 0.07 & 0.10 & 0.07 \\
slicing/array/lookup & 0.20 & 5.03 & 0.10 & 11.43 & 0.43 & 420.00 & 3.63 & 2.97 & 3.67 & 3.57 \\
slicing/array/dims & 0.20 & 0.07 & 0.10 & 0.47 & 0.10 & 24.00 & 0.03 & 0.10 & 0.07 & 0.07 \\
slicing/convolution/edgeDetect & 126.83 & 336.43 & 5.63 & 1068.73 & 9.93 & 27045.00 & 312.10 & 279.90 & 327.77 & 329.30 \\
slicing/convolution/emboss & 117.07 & 243.83 & 4.07 & 858.23 & 6.57 & 21464.00 & 267.20 & 238.67 & 304.50 & 290.70 \\
slicing/convolution/gaussian & 125.50 & 242.67 & 4.73 & 855.40 & 6.77 & 21464.00 & 271.80 & 238.80 & 288.93 & 300.27 \\
slicing/dict/create & 0.23 & 0.20 & 0.20 & 0.33 & 0.13 & 15.00 & 0.07 & 0.10 & 0.03 & 0.00 \\
slicing/dict/difference & 0.30 & 0.27 & 0.13 & 0.43 & 0.13 & 24.00 & 0.10 & 0.13 & 0.13 & 0.07 \\
slicing/dict/disjointUnion & 0.07 & 0.20 & 0.20 & 0.33 & 0.13 & 19.00 & 0.03 & 0.03 & 0.07 & 0.03 \\
slicing/dict/foldl & 0.27 & 1.57 & 0.17 & 3.63 & 0.30 & 110.00 & 1.27 & 1.67 & 1.37 & 1.30 \\
slicing/dict/intersectionWith & 0.20 & 0.43 & 0.20 & 0.77 & 0.27 & 47.00 & 0.13 & 0.43 & 0.20 & 0.20 \\
slicing/dict/fromRecord & 0.07 & 0.03 & 0.03 & 0.17 & 0.10 & 10.00 & 0.00 & 0.07 & 0.13 & 0.03 \\
slicing/dict/get & 0.13 & 0.33 & 0.17 & 0.40 & 0.17 & 29.00 & 0.07 & 0.13 & 0.13 & 0.10 \\
slicing/dict/map & 0.17 & 0.87 & 0.13 & 1.43 & 0.30 & 71.00 & 0.23 & 0.60 & 0.33 & 0.23 \\
slicing/divide & 0.10 & 0.07 & 0.13 & 0.10 & 0.10 & 7.00 & 0.00 & 0.03 & 0.00 & 0.10 \\
slicing/dtw/compute-dtw & 6245.07 & 75.67 & 6.27 & 231.30 & 26.23 & 6013.00 & 67.13 & 76.07 & 71.50 & 66.00 \\
slicing/dtw/average-series & 4933.70 & 107.53 & 1.60 & 243.17 & 15.70 & 6544.00 & 67.63 & 76.33 & 75.27 & 69.30 \\
slicing/filter & 0.43 & 0.47 & 0.17 & 1.73 & 0.33 & 54.00 & 0.50 & 0.60 & 0.70 & 0.70 \\
slicing/intersperse & 0.43 & 0.53 & 0.10 & 1.17 & 0.40 & 51.00 & 0.33 & 0.37 & 0.40 & 0.27 \\
slicing/intersperse & 0.40 & 0.53 & 0.20 & 1.37 & 0.60 & 51.00 & 0.40 & 0.30 & 0.47 & 0.30 \\
slicing/length & 0.27 & 0.27 & 0.27 & 1.27 & 0.63 & 48.00 & 0.20 & 0.27 & 0.33 & 0.27 \\
slicing/list-comp & 1.60 & 3.47 & 0.50 & 14.17 & 1.63 & 379.00 & 4.93 & 5.33 & 5.23 & 5.00 \\
slicing/list-comp & 1.63 & 3.53 & 0.47 & 13.93 & 1.53 & 379.00 & 4.77 & 4.87 & 4.53 & 4.73 \\
slicing/lookup & 0.77 & 0.67 & 0.27 & 2.30 & 0.90 & 96.00 & 0.67 & 0.57 & 0.63 & 0.57 \\
slicing/map & 0.17 & 0.23 & 0.13 & 0.57 & 0.30 & 29.00 & 0.17 & 0.17 & 0.23 & 0.23 \\
slicing/matrix-update & 0.37 & 22.27 & 0.40 & 48.07 & 1.40 & 1554.00 & 13.57 & 12.80 & 13.87 & 14.40 \\
slicing/multiply & 0.10 & 0.03 & 0.07 & 0.17 & 0.07 & 11.00 & 0.03 & 0.03 & 0.00 & 0.03 \\
slicing/nth & 0.07 & 0.47 & 0.10 & 0.83 & 0.17 & 38.00 & 0.20 & 0.17 & 0.27 & 0.27 \\
slicing/output-not-source & 0.07 & 0.00 & 0.07 & 0.20 & 0.07 & 9.00 & 0.00 & 0.03 & 0.07 & 0.00 \\
slicing/section-5-example & 2.13 & 3.53 & 0.37 & 14.30 & 1.73 & 421.00 & 4.63 & 4.27 & 4.73 & 4.67 \\
slicing/section-5-example & 1.37 & 4.27 & 0.40 & 14.63 & 1.30 & 421.00 & 5.03 & 4.60 & 4.77 & 4.97 \\
slicing/section-5-example & 1.80 & 3.83 & 0.40 & 14.43 & 1.40 & 421.00 & 4.73 & 4.37 & 4.57 & 4.83 \\
slicing/zeros & 0.20 & 0.13 & 0.13 & 0.43 & 0.17 & 23.00 & 0.10 & 0.10 & 0.17 & 0.10 \\
slicing/zeros & 0.13 & 0.10 & 0.10 & 0.53 & 0.13 & 23.00 & 0.07 & 0.07 & 0.20 & 0.07 \\
slicing/zipWith & 0.67 & 0.60 & 0.33 & 2.23 & 0.53 & 97.00 & 0.57 & 0.40 & 0.67 & 0.47 \\
graphics/background & 0.03 & 1.33 & 0.20 & 1.87 & 0.40 & 91.00 & 0.17 & 0.23 & 0.17 & 0.17 \\
graphics/grouped-bar-chart & 2.37 & 1144.60 & 0.90 & 706.60 & 546.87 & 2985.00 & 34.57 & 29.53 & 34.13 & 34.33 \\
graphics/line-chart & 1.30 & 490.37 & 1.10 & 593.77 & 214.47 & 4359.00 & 48.03 & 41.27 & 48.10 & 48.83 \\
graphics/stacked-bar-chart & 1.53 & 705.40 & 0.73 & 434.53 & 337.67 & 2045.00 & 19.30 & 16.13 & 19.10 & 18.93 \\
\bottomrule
\end{tabular}
\end{table}
}
